\documentclass[12pt]{article}
\usepackage[utf8]{inputenc}
\usepackage{indentfirst}
\usepackage{amsmath}

\usepackage{fancyhdr}
\fancyhead{}
\fancyfoot{}

\usepackage{color}

\usepackage[hidelinks,pdfauthor={Anamàs},pdftitle={Pla tècnic neteja comunitats}]{hyperref}

\usepackage{anysize}
\marginsize{3cm}{2.7cm}{2.7cm}{2.7cm}

\title{Pla tècnic}
\date{}

\begin{document}
\maketitle

\section{Càlcul de coordenades}

Molt ràpid.
El programa en CPP tarda al voltant de minuts en calcular les coordenades lat i lon de les 400 comunitats.

\section{Càlcul de distàncies}

Lent.
El programa en CPP calcula al voltant de 14-42 distàncies de parells de punts (origen, destí) en 1 minut.
Com que s'han de calcular 66.000 valors o distàncies entre les 400 comunitats (assumim una matriu simètrica, pel que no calculem $d(2,1)$ si ja hem calculat $d(1,2)$, ja que assumim que són iguals).

Es tarden 66.000/(14*60) = 79h = 3 dies i 7h en calcular tots els parells de distàncies a peu, metro, cotxe i bici entre tots els punts com a mínim.
Si la velocitat és màxima (42 distàncies per minut) es tardaria 1 dia i 2h en calcular-les totes.

\section{Descripció}

\subsection{Plantejament del problema}

Necessitem fer assignacions de treballadors amb comunitats i hores del dia.
Dit d'una altra manera, per a cada treballador volem calcular un itinerari durant cadascun dels dies de la setmana, i amb unes hores aproximades que minimitzin el temps de trajecte global en cadascun dels dies.

En concret, el nostre objectiu és minimitzar el temps de trajecte total, amb les següents limitacions:

\begin{enumerate}
	\item Cadascuna de les comunitats s'ha de netejar un nombre determinat de dies a la setmana, entre $1$ i $6$:
		\begin{enumerate}
			\item Si la comunitat es neteja $6$, $5$, o $3$ dies a la setmana, ha de ser en uns dies concrets (dl-ds, dl-dv, i dl,dc,dv respectivament), i per tant el nostre model no ha de prendre cap decisió al respecte.
			\item Si la comunitat es neteja $1$ dia a la setmana, aquest dia pot ser qualsevol feiner (dl, dt, dc, dj, dv) (i el nostre model haurà de decidir quin dia i hora assignar-li).
			\item Similarment, si la comunitat es neteja $4$ dies a la setmana, aquests hauran de ser feiners, i el model haurà d'assignar hores per cadascun dels 5 dies que triï.
			\item \label{item:2diessetmana} Si la comunitat es neteja $2$ dies a la setmana, només hi ha dues combinacions de dies que acceptem: o bé (dl,dj) o bé (dt,dv) (amb la corresponent tria d'horaris).
		\end{enumerate}

	\item Les claus de cadascuna de les comunitats es donaran a un únic treballador.
		És a dir, cada comunitat tindrà un treballador concret que la visitarà en tots els horaris que el model li assigni per netejar-la.
	\item Cada treballador pot treballar un màxim de 7h i 30m cada dia, incloent els desplaçaments.
	\item Cada treballador pot treballar un màxim de 37h i 30m cada setmana, incloent els desplaçaments.
	\item Festius? $\rightarrow$ TODO
\end{enumerate}

Les variables a optimitzar són les següents:

\begin{itemize}
	\item El temps de desplaçament total (de tots els treballadors en tots els dies).
	\item El nombre de treballadors necessaris per a cobrir totes les comunitats.
\end{itemize}

\subsection{Enfocament del problema}

La nostra modelització mirarà d'atacar el problema dividint-lo en tres parts principals:

\begin{enumerate}
	\item Distribuïm les comunitats en dies de la setmana de tal manera que es minimitzin tots els trajectes possibles en cada dia donat.
		Dit d'una altra manera, distribuïm les comunitats en clústers geogràfics en cadascun dels dies.
	\item Assignem un trajecte òptim a cada dia que cobreixi totes les comunitats cada dia en un clúster donat, i el dividim en parts de menys de 7h30m.
	\item Assignem treballadors a cada comunitat i trajecte per a cada dia.
\end{enumerate}

Havent completat aquests passos incorporarem altres consideracions (què fem amb els festius? Com podria adaptar-se el model a incidències i a baixes imprevistes?).

\section{Modelització: Assignació de comunitats per dia de la setmana}
\label{sec:dia_setmana}

Modelitzem aquesta part del problema amb un model d'optimització lineal.

Les variables del problema seran binàries (prendran els valors $0$ o $1$), i seran:

\begin{itemize}
	\item $x_{d,c}$: aquesta variable valdrà $1$ si la comunitat $c \in \{1,...,N\}$ es visita el dia $d \in \{1,...,6\}$ o $0$ altrament.
	\item $y_{d,c_1,c_2}$: aquesta variable equival a $x_{d,c_1} \wedge x_{d,c_2}$.
		Això vol dir que la variable $y_{d,c_1,c_2}$ val $1$ si en el dia $d$ es visiten tant $c_1$ com $c_2$, i $0$ altrament.
\end{itemize}

A més, tenim diversos paràmetres.
Denotem amb $D$ la matriu de distàncies òptimes entre comunitats, amb components $d_{c_1,c_2}$ (assumint que és simètrica).

Amb això, la funció objectiu és
\[F = \sum_{1 \leq c_1 \leq N} \sum_{c_1 < c_2 \leq N} d_{c_1,c_2} \sum_{d = 1}^6 y_{d,c_1,c_2} .\]

A més, caldrà aplicar les següents restriccions al sistema:

\begin{enumerate}
	\item $y_{d,c_1,c_2} = x_{d,c_1} \wedge x_{d, c_2}$.
		Per tal d'expressar-ho de forma lineal, ho hem de fer amb aquest conjunt de desigualtats per a cada $d,c_1,c_2$:
		\begin{align*}
			y_{d,c_1,c_2} &\geq x_{d,c_1} + x_{d,c_2} - 1, \\
			y_{d,c_1,c_2} &\leq x_{d,c_1}, \\
			y_{d,c_1,c_2} &\leq x_{d,c_2}.
		\end{align*}

	\item Si la comunitat $c$ es visita $n_c$ cops per setmana, llavors
		\[\sum_{d=1}^6 x_{d,c} = n_c .\]

	\item Si $n_c \neq 6$, $x_{6,c} = 0$.

	\item Si la comunitat $c$ es visita 3, 5 o 6 cops per setmana llavors els valors $x_{d,c}$ estan determinats (per exemple, $x_{1,c} = x_{3,c} = x_{5,c} = 1$ i $x_{2,c} = x_{4,c} = x_{6,c} = 0$ si la comunitat $c$ es visita 3 cops per setmana).

	\item Si $n_c = 2$, cal afegir les condicions
		\begin{align*}
			x_{3,c} = 0, \\
			x_{1,c} = x_{4,c}, \\
			x_{2,c} = x_{5,c},
		\end{align*}
		per tal de garantir la distribució de dies de (\ref{item:2diessetmana}).
\end{enumerate}

\section{Model: Classificació de les comunitats en clústers}

Farem servir una tècnica d'arbres de classificació i regressió (vegeu CART \cite{cart}).

El nostre espai de comunitats vindrà donat per les coordenades $A = (\mathrm{lat}, \mathrm{lon}, a_1, \dots, a_6)$, on $(\mathrm{lat}, \mathrm{lon})$ denoten la latitud i longitud on es troba la comunitat, i $a_i$ és una variable binària que indica si la comunitat és visitada en el dia $i$ (1) o no (0).
Fent servir la notació de la secció \ref{sec:dia_setmana}, $a_i = x_{i, A}$, que ens vindrà donada per la solució del pas anterior.
Les coordenades s'hauran de normalitzar (dividint per la desviació típica estàndard de cada magnitud).

A partir d'aquí hem de seleccionar la mètrica que ens funcioni millor pel nostre model, amb el criteri que

\begin{enumerate}
	\item Volem que dues comunitats siguin més properes si coincideixen en més dies (és a dir, volem que tendeixin a coincidir en el mateix clúster).
	\item Si dues comunitats no coincideixen en \textbf{cap} dia, llavors no ens importa la distància entre elles\footnote{Això és perquè si dues comunitats no coincideixen en cap dia, llavors el desplaçament entre l'una i l'altra no es produeix en hores de feina si les acaba fent el mateix treballador, i per tant no afecta al nostre model.}.
		Per tant, dues comunitats que no coincideixin en cap dia, en principi, poden acabar en el mateix clúster sense afegir cap cost addicional.
\end{enumerate}

Si $A = (\mathrm{lat}_A, \mathrm{lon}_A, a_1, \dots, a_6)$ i $B = (\mathrm{lat}_B, \mathrm{lon}_B, b_1, \dots, b_6)$, anomenem $\mathrm{coin}(A,B) = 6 - |a_1 - b_1| - \cdots - |a_6 - b_6|$.

La possible distància que estem considerant ara mateix seria:
\[\mathrm{discoin}(A,B) = \begin{cases} k \|(\mathrm{lat}_A, \mathrm{lon}_A) - (\mathrm{lat}_B, \mathrm{lon}_B)\| + |a_1 - b_1| + \cdots + |a_6 - b_6| & \text{si } \mathrm{coin}(A,B) > 0, \\ 0 & \text{si } \mathrm{coin}(A,B) = 0, \end{cases}\]

	on $k > 0$ és un pes que associem a la distància geogràfica i ens permetria calibrar quanta importància li volem donar en el nostre model. 

\begin{thebibliography}{9}
	\bibitem{cart} Breiman, L., Friedman, J., Olshen, R. A., Stone, C. J. Classification and regression trees. Routledge, 1984.
\end{thebibliography}
\end{document}
